% Originally: content/week-08.tex, \subsection{Fubini's theorem} --- promoted to a section

\section{Fubini's theorem}
\label{sec:fubini}

\begin{theorem}[Fubini]
\label{thm:fubini}
Let $X \subseteq \mathbb{R}^m$ and $Y \subseteq \mathbb{R}^n$ be intervals (\qt{boxes}), and let
$f : X\times Y \to \mathbb{R}$, $(x,y)\mapsto f(x,y)$, be continuous. Then
\[\int_{X\times Y} f(x,y)\,dx\,dy
  = \int_X\left(\int_Y f(x,y)\,dy\right)dx
  = \int_Y\left(\int_X f(x,y)\,dx\right)dy,\]
where we look at the expression in brackets as a function, e.g.\
$x \mapsto \int_Y f(x,y)\,dy$.
\end{theorem}

% Source: Corsin Nick/Class Notes/Week 8.pdf, p. 8
\begin{example}[Fubini on the unit square]
\label{ex:fubini_unit_square}
\[
\begin{aligned}
\int_{[0,1]^2} x^{y+1}\,dx\,dy
  &= \int_0^1 dy \int_0^1 dx\ x^{y+1}
   = \int_0^1 dy\ \left[\frac{x^{y+2}}{y+2}\right]_0^1 \\
  &= \int_0^1 dy\ \frac{1}{y+2}
   = \log|y+2|\,\Big|_0^1
   = \log\!\left(\tfrac32\right).
\end{aligned}
\]
\end{example}

\begin{ainote}
Corsin writes the first line as $\int_0^1 dx \int_0^1 dy$, but the antiderivative taken on the
next line is in $x$ --- it produces $x^{y+2}/(y+2)$, evaluated from $0$ to $1$. The differentials
are therefore swapped relative to the computation actually performed; the order is corrected
above and the result is unaffected.

The example is well chosen: integrating $x^{y+1}$ in $y$ first would mean integrating an
exponential in the exponent, which is far worse. Fubini's real use is that you may pick the
easier order, and here only one of the two is comfortable.
\end{ainote}

% Supplement: Sascha Brack/Class Notes/Week_08_Notes_Friday.pdf, pp. 4--5
\begin{example}[Changing the order of integration]
\label{ex:change_order_integration}
Consider the iterated integral $\int_0^2 \int_{y^2}^{2\sqrt{2y}} f(x,y) \,dx\,dy$. Let us rewrite it with the integration order reversed (first $y$, then $x$).

The region is defined by $0 \leq y \leq 2$ and $y^2 \leq x \leq 2\sqrt{2y}$.
Notice that for $y \in [0,2]$, the lower bound is $x \geq 0$ and the upper bound is $x \leq 2\sqrt{4} = 4$, so $x$ ranges from $0$ to $4$.
We now solve the inequalities for $y$:
$x \leq 2\sqrt{2y} \implies x^2 \leq 8y \implies y \geq \frac{1}{8}x^2$.
$x \geq y^2 \implies y \leq \sqrt{x}$ (since $x,y \geq 0$).
Thus, the new bounds are $0 \leq x \leq 4$ and $\frac{1}{8}x^2 \leq y \leq \sqrt{x}$. The reversed integral is:
\[\int_0^4 \int_{\frac{1}{8}x^2}^{\sqrt{x}} f(x,y) \,dy\,dx.\]
\end{example}

\begin{example}[Surface parametrization and Fubini]
\label{ex:surface_parametrization_integral}
Let $f(u,v) := (u\cos v, u\sin v, u^2)\transp \in \mathbb{R}^3$ for $(u,v) \in \mathbb{R}^2$. Write the third coordinate $z(u,v)$ as a function of the first two, $\tilde{z}(x,y)$, and compute $\int_0^y \int_0^x \tilde{z}(x',y') \,dx'\,dy'$.

We have $x(u,v) = u\cos v$ and $y(u,v) = u\sin v$. Then $x^2 + y^2 = u^2\cos^2 v + u^2\sin^2 v = u^2 = z(u,v)$.
So we define $\tilde{z}(x,y) := x^2 + y^2$.
The integral is:
\begin{align}
\int_0^y \int_0^x (x'^2 + y'^2) \,dx'\,dy'
&= \int_0^y \left[ \frac{x'^3}{3} + x'y'^2 \right]_{x'=0}^{x'=x} \,dy' \nonumber \\
&= \int_0^y \left( \frac{x^3}{3} + x y'^2 \right) \,dy' \nonumber \\
&= \left[ \frac{x^3 y'}{3} + x \frac{y'^3}{3} \right]_{y'=0}^{y'=y} \nonumber \\
&= \frac{xy}{3}(x^2 + y^2). \label{eq:surface_integral_result}
\end{align}
\end{example}

% Source: Corsin Nick/Class Notes/Week 8.pdf, pp. 8--10
\begin{example}[A domain that Fubini alone cannot handle]
\label{ex:change_of_variables_domain}
Consider
\[\int_A \frac{y}{\sqrt{x}}\,dy\,dx, \qquad
A := \left\{(x,y) \in \mathbb{R}^2 \ :\ x,y>0,\ \ 1 \leq \frac{y}{\sqrt x} \leq 2,\ \ 1 \leq xy \leq 2\right\}.\]
This is not easily integrated with Fubini alone. But we can apply a transformation of variables
to simplify the integration domain $A$: take
\[\Phi : U \to V, \qquad \begin{pmatrix}x\\y\end{pmatrix} \mapsto
  \begin{pmatrix}u(x,y)\\v(x,y)\end{pmatrix} = \begin{pmatrix}\tfrac{y}{\sqrt x}\\[2pt] xy\end{pmatrix},\]
with $U$ and $V$ still to be determined --- they must be open.

\textbf{Check injectivity.} Suppose, for $x,y,x',y' > 0$,
\[\text{\textbf{(1)}}\ \ \frac{y}{\sqrt x} = \frac{y'}{\sqrt{x'}},
\qquad\qquad \text{\textbf{(2)}}\ \ xy = x'y'.\]
From \textbf{(2)}, $\dfrac{x}{x'} = \dfrac{y'}{y}$. Inserting into \textbf{(1)},
\[\frac{1}{\sqrt x} = \frac{y'}{y}\cdot\frac{1}{\sqrt{x'}} = \frac{x}{x'}\cdot\frac{1}{\sqrt{x'}}
\quad\Longrightarrow\quad x^{3/2} = (x')^{3/2} \quad\Longrightarrow\quad x = x',\]
and consequently $\dfrac{y'}{y} = \dfrac{x}{x'} = 1$. So for $U := (0,\infty)^2$ and
$V := \Phi(U)$, the map $\Phi$ is a bijection between open sets.

\textbf{Functional determinant.} We calculate $|\det \Jac\Phi|$; if it is strictly positive we
may proceed, since $\Phi$ is then a local diffeomorphism as well as a bijection, hence a
diffeomorphism. Now
\[\Jac\Phi(x,y) = \begin{pmatrix} -\dfrac{y}{2\sqrt{x}^{\,3}} & \dfrac{1}{\sqrt x} \\[2ex] y & x \end{pmatrix},\]
so that
\[\big|\det \Jac\Phi(x,y)\big| = \left|-\frac32\,\frac{y}{\sqrt x}\right| = \frac32\,\frac{y}{\sqrt x} = \frac32\,u(x,y).\]
So $\Phi$ is a diffeomorphism and $du\,dv = \tfrac32 u\,dx\,dy$. Furthermore, by design,
\[\Phi(A) = (1,2)^2.\]
So with the change of variables (\cref{thm:change_of_variables}), followed by Fubini
(\cref{thm:fubini}),
\[\int_A \frac{y}{\sqrt x}\,dy\,dx
  = \int_{(1,2)^2} u\,\frac{1}{\tfrac32 u}\,du\,dv
  = \frac23\int_1^2 dv \int_1^2 du
  = \frac23.\]
\end{example}

\begin{ainote}
The design principle is worth stating, because it is the reverse of how substitution is usually
taught. Here $\Phi$ was not chosen to simplify the \emph{integrand} but to straighten the
\emph{domain}: the two constraints defining $A$ are literally $1 \leq u \leq 2$ and
$1 \leq v \leq 2$, so $\Phi$ turns an awkward curvilinear region into a square, and only then
does Fubini become available. That the integrand $y/\sqrt x$ is itself $u$, and cancels against
most of the Jacobian factor, is a bonus of the same choice.
\end{ainote}

% Sonnet 5 (Medium)
We now want to look at a powerful new integration technique using ODEs.
